\documentclass[11pt]{article}

\usepackage[margin=1in]{geometry}
\usepackage{amsmath,amssymb,amsthm,mathtools,microtype}
\usepackage[T1]{fontenc}
\usepackage[hidelinks]{hyperref}

\newtheorem{theorem}{Theorem}[section]
\newtheorem{proposition}[theorem]{Proposition}
\newtheorem{lemma}[theorem]{Lemma}
\newtheorem{corollary}[theorem]{Corollary}
\newtheorem{remark}[theorem]{Remark}

\newcommand{\N}{\mathbb N}
\newcommand{\Q}{\mathbb Q}
\newcommand{\Z}{\mathbb Z}
\newcommand{\Newt}{\operatorname{Newt}}
\newcommand{\Supp}{\operatorname{Supp}}
\newcommand{\CT}{\operatorname{CT}}
\newcommand{\wt}{\operatorname{wt}}

\title{The Generalized Vanishing Conjecture:\\
The Two-Variable Theorem and the First Failing Dimension}
\author{Roy van Rijn\\{\small Independent researcher}}
\date{4 August 2026}
\hypersetup{
  pdftitle={The Generalized Vanishing Conjecture: The Two-Variable Theorem and the First Failing Dimension},
  pdfauthor={Roy van Rijn},
  pdfsubject={The exact dimensional boundary of the Generalized Vanishing Conjecture},
  pdfkeywords={Generalized Vanishing Conjecture, constant-coefficient differential operators,
    Newton polygons, Hall's theorem, Gaussian moments, counterexample}
}

\begin{document}
\maketitle

\begin{abstract}
We determine the exact dimensional range of the Generalized Vanishing
Conjecture for constant-coefficient differential operators over a field of
characteristic zero.  In two variables we prove that
\[
 \Lambda^m(P^m)=0\quad(m\geq1)
 \qquad\Longrightarrow\qquad
 \Lambda^m(QP^m)=0\quad(m\gg0)
\]
for every fixed polynomial $Q$.  The proof follows two moving faces.  Hall's
marriage theorem puts the lowest operator face strictly to the right of the
highest polynomial face near ordinary degree.  A prime-dilation argument
shows that, while the two faces have different weights, their horizontal
Newton intervals cannot overlap.  Since the supports are finite, the two
piecewise-linear weight envelopes must eventually meet.  Their first meeting
is an unequal-weight common threshold, where a fixed multiplier cannot repair
the resulting linear derivative deficit.  The only deep external theorem is
the constant-term theorem of Duistermaat and van der Kallen.  We then give a
homogeneous
three-variable counterexample.  With
$\rho=t^2+xy$, $A=\rho+x^2$,
$C=y\rho^2-2xt^2\rho-x^3t^2$, and
$\Delta=4\partial_x\partial_y+\partial_t^2$, the pair
\[
 \Lambda=\Delta^6,\qquad P=AC^2
\]
satisfies $\Lambda^m(P^m)=0$ while
$\Lambda^m(x^2P^m)\ne0$ for every $m\geq1$.  The construction is a
homogeneous endpoint-contact lift of Long's three-dimensional Gaussian
counterexample.  It is the smallest member of an explicit cusp-profile
family with arbitrary winding, profile, and radial suspension; in
particular every $\Delta^k$, $k\geq6$, has such a failure.  Consequently
unrestricted, and already homogeneous, GVC holds in dimension $n$ exactly
when $n\leq2$.
\end{abstract}

\section{Introduction}

For a polynomial $\lambda(X,Y)$, write
\[
 \Lambda=\lambda(\partial_x,\partial_y).
\]
The Generalized Vanishing Conjecture asks whether, for all polynomials
$P,Q$ over a characteristic-zero field,
\begin{equation}\label{eq:gvc}
 \Lambda^m(P^m)=0\quad\text{for every }m\geq1
 \quad\Longrightarrow\quad
 \Lambda^m(QP^m)=0\quad\text{for all sufficiently large }m.
\end{equation}
It grew out of the Vanishing Conjecture and its relation to the Jacobian
Conjecture; see Zhao~\cite{Zhao} and van den Essen, Willems, and
Zhao~\cite{EWZ}.

There were several particularly important precursors.  De Bondt proved GVC
for every product of linear differential forms, and therefore for every
homogeneous operator in two variables~\cite{deBondt}.  Feng and Sun proved
the nonhomogeneous binary case
\[
 \Lambda=(\partial_x-\Phi(\partial_y))\partial_y
\]
for an arbitrary one-variable polynomial $\Phi$~\cite{FengSun}.  Wilson's
independent July 2026 preprint proves the two-variable Gaussian Moments
Conjecture by a related Newton-face and prime-isolation method
~\cite{WilsonGMC}.  Wilson treats the Gaussian functional rather than an
arbitrary constant-coefficient operator.  In the present argument, Hall
localization and separation on shifted output rays are the additional steps
that pass from that philosophy to unrestricted binary GVC.

Also in 2026, Long found a five-term quartic in three real Gaussian variables for which
all pure Gaussian moments vanish while one fixed mixed moment equals $m!$
at every power~\cite{LongGMC}.  Long's example disproves the Gaussian
Moments Conjecture in every dimension at least three, but it is not itself a
constant-coefficient GVC counterexample.

This paper closes the gap on both sides.  First we prove \eqref{eq:gvc} for
arbitrary, possibly nonhomogeneous, operators in two variables.  We then
lift the endpoint-coefficient mechanism behind Long's example to a
homogeneous polynomial and a homogeneous differential operator in three
variables.  The two results meet at the exact boundary.

\begin{theorem}[Binary Generalized Vanishing Conjecture]\label{thm:main}
Let $k$ be a field of characteristic zero.  Every constant-coefficient
differential operator $\Lambda$ on $k[x,y]$ satisfies
\eqref{eq:gvc}.
\end{theorem}

The global proof is organized as follows.  Give both
the operator variables and the polynomial variables weight $(s,1)$.  Let
$L(s)$ be the least operator weight and $U(s)$ the greatest polynomial
weight.  Hall's theorem orders the corresponding faces just to the right of
$s=1$.  Prime-dilated endpoint separation prevents the faces from passing
one another while $U(s)-L(s)>0$.  Finite support then forces this difference
to reach zero.  At the first zero all operator terms lie above one threshold
and all polynomial terms lie below it, which gives eventual mixed
vanishing.

For orientation, the load-bearing dependencies in the binary argument are
\[
\begin{array}{c}
 \text{Duistermaat--van der Kallen}+\text{split-symbol polarization}
   \longrightarrow \text{Proposition~\ref{prop:cutoff}},\\[2pt]
 \text{Duistermaat--van der Kallen}+\text{Hall localization}
   \longrightarrow \text{Lemma~\ref{lem:hall}},\\[2pt]
 \text{rational-point exposure}+\text{prime isolation}
   \longrightarrow \boxed{\text{Proposition~\ref{prop:shifted}}}
   \longrightarrow \text{Corollary~\ref{cor:horizontal}},\\[2pt]
 \text{Lemma~\ref{lem:hall}}+\text{Corollary~\ref{cor:horizontal}}
   \longrightarrow \text{first common threshold}
   \longrightarrow \text{Proposition~\ref{prop:threshold}}
   \longrightarrow \text{Theorem~\ref{thm:main}}.
\end{array}
\]
Thus Proposition~\ref{prop:shifted} is the single arithmetic bridge from
the local Hall gap to the global envelope argument.

\section{Two standard reductions}

We use the Newton polytope of a Laurent polynomial $F$,
\[
 \Newt(F)=\operatorname{conv}(\Supp(F)).
\]
The only substantial external input is the following theorem.

\begin{theorem}[Duistermaat--van der Kallen~\cite{DK}]\label{thm:dk}
Let $F$ be a Laurent polynomial over a characteristic-zero field.  If
$\CT(F^m)=0$ for every $m\geq1$, then
$0\notin\Newt(F)$.
\end{theorem}

We will also use the following immediate consequence.

\begin{lemma}[Exposing a rational point]\label{lem:expose}
Let $F$ be a nonzero Laurent polynomial and let
$\rho\in\Newt(F)\cap\Q^n$.  Then there is an integer $a\geq1$ such that
\[
 [z^{a\rho}]F^a\ne0.
\]
For finitely many polynomials and rational points, the same $a$ may be
chosen for all of them.
\end{lemma}

\begin{proof}
Choose $N$ with $N\rho\in\Z^n$ and apply
Theorem~\ref{thm:dk} contrapositively to
$z^{-N\rho}F^N$.  To use one exponent for several polynomials, place them
in disjoint sets of Laurent variables, multiply them, and apply the same
argument to the product.
\end{proof}

We next dispose of the range in which the degree of $P$ is no larger than
the first positive order of $\Lambda$.  This familiar split-symbol argument
is included so that the main theorem has no hidden degree case.

\begin{proposition}[Binary degree cutoff]\label{prop:cutoff}
In two variables, suppose that $\Lambda$ has no constant term, and let $r$
be the least
positive degree appearing in its symbol.  If $\deg P\leq r$, then
\eqref{eq:gvc} holds.
\end{proposition}

\begin{proof}
Put $d=\deg P$.  If $d<r$, then
\[
 \operatorname{ord}(\Lambda^m)\geq rm>\deg(QP^m)
\]
for all large $m$.

Assume $d=r$.  Let $\lambda_r$ and $P_r$ be the lowest symbol and the top
homogeneous part of $P$.  Over an algebraic closure,
\[
 \lambda_r=\prod_{i=1}^r D_{v_i}.
\]
With $Vt=\sum_i t_iv_i$, complete polarization gives
\[
 \lambda_r(\partial)^m(P_r^m)
 =(m!)^r\CT_t\left(
       \frac{P_r(Vt)}{t_1\cdots t_r}
                  \right)^m.
\]
The highest-degree part of the pure premise is zero, so if $P_r(Vt)\ne0$,
Theorem~\ref{thm:dk} supplies an integral linear functional $h$ which is
strictly positive on the support of the displayed Laurent polynomial.  We
first treat this case.

Write $\Lambda=\sum_{j=r}^R\Lambda_j$ by homogeneous order and expand
$\Lambda^m$.  If $n_j$ factors equal $\Lambda_j$, put
\[
 s=\sum_{j>r}n_j,
 \qquad e=\sum_{j>r}(j-r)n_j.
\]
The corresponding term is $\Lambda_r^{m-s}B$, where the order of the fixed
operator $B$ is $rs+e$.  Its total order is $rm+e$, so it can act
nontrivially on $QP^m$ only when $e\leq\deg Q$.  Consequently $s$, $e$, and
$B$ range over a finite set independent of $m$.

Expand the bounded operator $B$ by Leibniz.  Each resulting term is
$\Lambda_r^{m-s}$ applied to a product of one fixed derivative of $Q$ and
$m$ copies of $P$, only a bounded number of which have been differentiated.
Polarization asks for the coefficient of
$t_1^{m-s}\cdots t_r^{m-s}$.  Its total $t$-degree is $r(m-s)$.  Since a
non-top choice from an undifferentiated copy of $P$ loses at least one unit
of total $t$-degree, while the translated $Q$ factor and all differentiated
copies have bounded degree, at least $m-C$ copies must contribute a monomial
of $P_r(Vt)$.

Put $H_0=h(1,\ldots,1)$.  The separator says that every such top monomial
$t^\gamma$ satisfies $h(\gamma)\geq H_0+1$.  All remaining factors have
total $h$-weight bounded below by a constant.  Hence every possible product
has $h$-weight at least
\[
 (m-C)(H_0+1)-C',
\]
whereas the requested diagonal exponent has weight $(m-s)H_0$.  The former
is larger for all sufficiently large $m$, so the coefficient is zero.  This
kills every term of $\Lambda^m(QP^m)$.

It remains to treat $P_r(Vt)=0$.  Then $P(z+Vt)$ has total $t$-degree at
most $r-1$.
Derivatives do not increase this degree, so after the bounded insertion the
whole translated product has $t$-degree at most $(r-1)m+C$.  The diagonal
target for $\lambda_r^{m-s}$ has degree $r(m-s)$, which is larger for all
large $m$.  The argument is unchanged by scalar extension, so it descends
to $k$.
\end{proof}

If $\lambda$ has a nonzero constant term $c$, the top ordinary-degree part
of $\Lambda^m(P^m)$ is $c^mP_d^m$.  Thus the pure premise forces $P=0$.
The zero operator is also trivial.  From now on we may therefore assume
that $r\geq1$ and $d=\deg P>r$.

\section{The Hall starting point}

Write $\lambda_r$ for the degree-$r$ part of the symbol and $P_d$ for the
degree-$d$ part of $P$.  The highest ordinary-degree part of the pure
identity is
\begin{equation}\label{eq:top-pure}
 \lambda_r(\partial)^m(P_d^m)=0\qquad(m\geq1).
\end{equation}

\begin{lemma}[Binary Hall localization]\label{lem:hall}
After extending the field and making a linear change of coordinates, there
is an integer $e$, $1\leq e\leq r$, such that
\begin{equation}\label{eq:hall-normal}
 \lambda_r=X^eC_{r-e}(X,Y),\qquad C_{r-e}(0,1)\ne0,
\end{equation}
and
\begin{equation}\label{eq:hall-poly}
 P_d=y^{d-e+1}D_{e-1}(x,y).
\end{equation}
In particular, the least $X$-exponent of $\lambda_r$ is $e$, while the
greatest $x$-exponent of $P_d$ is at most $e-1$.
\end{lemma}

\begin{proof}
Factor
\[
 \lambda_r=\prod_{i=1}^r D_{v_i},
 \qquad P_d=\prod_{j=1}^d L_j
\]
over an algebraic closure.  Introduce $t_1,\ldots,t_r$ and a generic
translation point $z$.  Polarization turns \eqref{eq:top-pure} into
\[
 \CT_t\left(
   \frac{P_d(z+\sum_i t_iv_i)}{t_1\cdots t_r}
          \right)^m=0
 \qquad(m\geq1).
\]
By Theorem~\ref{thm:dk}, the vector $(1,\ldots,1)$ does not belong to the
Newton polytope of $P_d(z+\sum_i t_iv_i)$.

For each derivative copy $i$, join $i$ to the factor $L_j$ when
$L_j(v_i)\ne0$.  If Hall's inequalities held, there would be a matching
from all derivative copies to distinct polynomial factors.  Choosing the
matched $t_i$ term in those factors and the constant term in the others
would put $(1,\ldots,1)$ in the Minkowski sum of their Newton polytopes,
which is the Newton polytope of the product.  Hall's theorem therefore gives
a set $I$ of derivative copies with fewer than $|I|$ neighboring polynomial
factors.

The vectors belonging to $I$ must all be parallel.  Indeed, two
nonparallel vectors cannot both be annihilated by a nonzero binary linear
form, so such a pair would see all $d>r$ polynomial factors and could not
be Hall-deficient.  Let $\mu=|I|$, and let $e\geq\mu$ be the full
multiplicity of this direction in $\lambda_r$.  At least $d-\mu+1$, and
hence at least $d-e+1$, factors of $P_d$ annihilate the direction.  Choosing
it as the $x$-axis gives \eqref{eq:hall-normal} and
\eqref{eq:hall-poly}.
\end{proof}

The last sentence of the lemma is the initial horizontal gap.  The rest of
the proof shows that this gap cannot close while the two extremal faces have
different weights.

We record separately the factorial calculation used in the prime-dilation
argument.  For a prime ideal $\mathfrak p$, normalize
$v_{\mathfrak p}$ by $v_{\mathfrak p}(\mathfrak p)=1$.

\begin{lemma}[Factorial valuations at an unramified prime]
\label{lem:factorial-valuations}
Let $K$ be a number field, let $p$ be a rational prime unramified in $K$,
and let $\mathfrak p\mid p$.  For $\eta,\gamma\in\N^2$, put
\[
 F_p(\gamma,\eta)=
 \prod_{i=x,y}\frac{(\gamma_i+p\eta_i)!}{(p\eta_i)!}.
\]
Then
\begin{equation}\label{eq:factorial-lower}
 v_{\mathfrak p}\bigl(F_p(\gamma,\eta)\bigr)
 \geq
 \left\lfloor\frac{\gamma_x}{p}\right\rfloor+
 \left\lfloor\frac{\gamma_y}{p}\right\rfloor.
\end{equation}
If $\gamma=p\beta$ and
$p>\max_i(\beta_i+\eta_i)$, then
\begin{equation}\label{eq:factorial-exact}
 v_{\mathfrak p}\bigl(F_p(p\beta,\eta)\bigr)
 =\beta_x+\beta_y.
\end{equation}
\end{lemma}

\begin{proof}
Since $p$ is unramified,
$v_{\mathfrak p}(n)=v_p(n)$ for every nonzero rational integer $n$.
Legendre's formula therefore gives, in each coordinate,
\[
 v_{\mathfrak p}\!\left(
  \frac{(\gamma_i+p\eta_i)!}{(p\eta_i)!}
 \right)
 =\sum_{j\geq1}\left(
  \left\lfloor\frac{\gamma_i+p\eta_i}{p^j}\right\rfloor-
  \left\lfloor\frac{p\eta_i}{p^j}\right\rfloor
 \right).
\]
The $j=1$ summand is $\lfloor\gamma_i/p\rfloor$, and all later
summands are nonnegative, proving \eqref{eq:factorial-lower}.  If
$\gamma_i=p\beta_i$ and $\beta_i+\eta_i<p$, the $j=1$ summand is
$\beta_i$, while every summand with $j\geq2$ is zero.  Summing over
$i=x,y$ proves \eqref{eq:factorial-exact}.
\end{proof}

\section{Prime-dilated endpoint separation}

Fix a positive weight $w=(u,v)$ with $u\ne v$.  We give $X,Y$ the same
weights as $x,y$.  A polynomial supported on one affine line
$u a+v b=W$ will be called $w$-homogeneous of weight $W$.

\begin{proposition}[Shifted-ray separation]\label{prop:shifted}
Let $A(X,Y)$ and $B(x,y)$ be nonzero $w$-homogeneous polynomials of
weights $W_A$ and
\[
 W_B=W_A+w(\delta),
 \qquad \delta\in\Q_{\geq0}^2.
\]
If
\begin{equation}\label{eq:shifted-vanishing}
 [x^{m\delta_x}y^{m\delta_y}]
 A(\partial)^m(B^m)=0
\end{equation}
at every integral scale after the denominator of $\delta$ is cleared, then
\begin{equation}\label{eq:shifted-disjoint}
 \Newt(A)\cap(\Newt(B)-\delta)=\varnothing.
\end{equation}
\end{proposition}

\begin{proof}
Suppose the intersection
\[
 I=\Newt(A)\cap\bigl(\Newt(B)-\delta\bigr)
\]
is nonempty.  Both Newton sets in this intersection lie on the affine line
of weight $W_A$.  Because $u\ne v$, ordinary degree is strictly monotone
on that line.  The compact segment $I$ therefore has a unique point
$\rho$ of least ordinary degree.

\emph{Exposure and specialization.}
Apply Lemma~\ref{lem:expose} in two independent sets of variables, at
$\rho$ for $A$ and at $\rho+\delta$ for $B$.  There is an integer $k\geq1$
which clears both lattice denominators and for which
\[
 c=[X^{k\rho}]A^k\ne0,
 \qquad b=[x^{k(\rho+\delta)}]B^k\ne0.
\]
Put
\[
 A_0=A^k,\qquad B_0=B^k,
 \qquad \alpha=k\rho,\qquad \eta=k\delta,
\]
so that $\alpha,\eta\in\N^2$.  Let $R$ be the finitely generated
$\Q$-subalgebra of the coefficient field generated by the coefficients of
$A,B$ and by $b^{-1},c^{-1}$.  It is a domain.  A maximal ideal of
$R$ has number-field residue field by Zariski's lemma.  Specializing to
that residue field preserves $b,c\ne0$.  It also preserves every identity
in \eqref{eq:shifted-vanishing}: each requested coefficient is an
integer-coefficient polynomial in the finitely many coefficients of $A$
and $B$, and a zero element of $R$ remains zero after specialization.
We may consequently replace the coefficient field by a number field $K$
without losing either the vanishing identities or the two exposed endpoint
coefficients.

\emph{Choice of the prime.}
Let
\[
 M=\max\{\nu_i:\nu\in\Supp(A_0)\cup\Supp(B_0),\ i=x,y\}.
\]
Exclude the following finite set of rational primes:
\begin{itemize}
 \item primes below a denominator of a coefficient of $A_0$ or $B_0$;
 \item primes ramified in $K$;
 \item primes below a prime ideal at which $b$ or $c$ is not a unit; and
 \item primes $p\leq M$.
\end{itemize}
Choose a rational prime $p$ outside this set and a prime ideal
$\mathfrak p\mid p$.  All coefficients of $A_0,B_0$ are then
$\mathfrak p$-integral, $b,c$ are $\mathfrak p$-adic units, $p$ is
unramified, and $p$ exceeds every coordinate needed below.

In the residue field of $\mathfrak p$, Frobenius gives
\begin{align}
 A_0^p&\equiv\sum_{\nu\in\Supp(A_0)}
   ([X^\nu]A_0)^pX^{p\nu}\pmod{\mathfrak p},
 \label{eq:frobenius-a}\\
 B_0^p&\equiv\sum_{\nu\in\Supp(B_0)}
   ([x^\nu]B_0)^px^{p\nu}\pmod{\mathfrak p}.
 \label{eq:frobenius-b}
\end{align}
In particular,
\[
 [X^\gamma]A_0^p\equiv0\pmod{\mathfrak p}
 \quad\text{unless }\gamma\in p\Z^2,
\]
and the analogous statement holds for $B_0^p$.  Since $p\eta\in p\Z^2$,
$[x^{\gamma+p\eta}]B_0^p$ also vanishes modulo $\mathfrak p$ unless
$\gamma\in p\Z^2$.

\emph{The coefficient and its valuations.}
The coefficient of $x^{p\eta_x}y^{p\eta_y}$ in
$A_0(\partial)^p(B_0^p)$ is
\begin{equation}\label{eq:coefficient-sum}
 \sum_\gamma
 F_p(\gamma,\eta)
 [X^\gamma]A_0^p\,
 [x^{\gamma+p\eta}]B_0^p.
\end{equation}
Every contributing $\gamma$ satisfies
\[
 \frac{\gamma}{p}\in
 \Newt(A_0)\cap\bigl(\Newt(B_0)-\eta\bigr)=kI.
\]
Let $s=\alpha_x+\alpha_y$, the least ordinary degree on $kI$.

First suppose $\gamma\notin p\Z^2$.  Write
$\gamma_i=pn_i+r_i$, where $0\leq r_i<p$.  Minimality of $s$ gives
\[
 n_x+n_y+\frac{r_x+r_y}{p}
 =\frac{\gamma_x+\gamma_y}{p}\geq s.
\]
Since $(r_x+r_y)/p<2$ and $n_x+n_y$ is an integer, this implies
\begin{equation}\label{eq:floor-bound}
 \left\lfloor\frac{\gamma_x}{p}\right\rfloor+
 \left\lfloor\frac{\gamma_y}{p}\right\rfloor
 =n_x+n_y\geq s-1.
\end{equation}
Lemma~\ref{lem:factorial-valuations} and
\eqref{eq:frobenius-a}--\eqref{eq:frobenius-b} now show that this summand
has $\mathfrak p$-adic valuation at least $(s-1)+1+1=s+1$.

Now suppose $\gamma=p\beta$.  Then $\beta\in kI\cap\Z^2$.  Moreover,
$\beta\in\Newt(A_0)$ and $\beta+\eta\in\Newt(B_0)$, so the definition of
$M$ and the choice $p>M$ give
$p>\max_i(\beta_i+\eta_i)$.  Hence
\begin{equation}\label{eq:divisible-factorial}
 v_{\mathfrak p}\bigl(F_p(p\beta,\eta)\bigr)
 =\beta_x+\beta_y
\end{equation}
by Lemma~\ref{lem:factorial-valuations}.  The unique point of $kI$ of
ordinary degree $s$ is $\alpha$.  Thus every integral
$\beta\ne\alpha$ has $\beta_x+\beta_y\geq s+1$.

At $\beta=\alpha$, equations \eqref{eq:frobenius-a} and
\eqref{eq:frobenius-b} give
\[
 [X^{p\alpha}]A_0^p\equiv c^p\not\equiv0\pmod{\mathfrak p},
 \qquad
 [x^{p(\alpha+\eta)}]B_0^p\equiv b^p\not\equiv0
 \pmod{\mathfrak p}.
\]
Together with \eqref{eq:divisible-factorial}, this says that the
$\gamma=p\alpha$ summand has valuation exactly $s$, whereas every other
summand has valuation at least $s+1$.  It is therefore the unique summand
of least $\mathfrak p$-adic valuation, so \eqref{eq:coefficient-sum} is
nonzero.  But this is precisely the specialization of the coefficient in
\eqref{eq:shifted-vanishing} at the admissible scale $kp$, a contradiction.
\end{proof}

\begin{remark}[Relation with Wilson's prime-isolation argument]
The arithmetic mechanism in Proposition~\ref{prop:shifted} closely parallels
the face-isolation argument of Wilson~\cite{WilsonGMC}.  Both arguments
first pass from the original coefficient domain to a number field and then
evaluate a moment at an index equal to a fixed integer times a sufficiently
large good prime: $m=qp$ in Wilson's notation and $m=kp$ above.
Frobenius separates a distinguished face or endpoint contribution from the
remaining terms.  A factorial-valuation calculation shows that the
distinguished contribution has its predicted exact $p$-adic valuation,
whereas every other contribution has valuation at least one greater.
Together with the Duistermaat--van der Kallen theorem, this turns the
vanishing of all moments into a Newton-support separation statement.

The geometric settings and the order of these ingredients are different.
Wilson isolates a lower face for the one-radial Gaussian functional.  Here
the distinguished term is the least-ordinary-degree endpoint $\alpha$ of a
shifted intersection, the relevant factorial factor is the two-coordinate
quotient $F_p(\gamma,\eta)$, and Duistermaat--van der Kallen enters through
Lemma~\ref{lem:expose} before the prime-dilation step.  Thus
Proposition~\ref{prop:shifted} is a shifted-ray analogue of Wilson's
prime-isolation mechanism, rather than a direct specialization of it.
The subsequent passage to arbitrary binary constant-coefficient operators
also uses Hall localization and the Newton-envelope argument.
\end{remark}

\begin{corollary}[Horizontal separation]\label{cor:horizontal}
Let $A$ and $B$ be nonzero $w$-homogeneous polynomials with
$W_A<W_B$.  If
\[
 A(\partial)^m(B^m)=0\qquad(m\geq1),
\]
then the projections of $\Newt(A)$ and $\Newt(B)$ to the $x$-axis are
disjoint.
\end{corollary}

\begin{proof}
If the projected intervals overlap, choose a rational number $c$ in their
intersection.  There are unique points
\[
 \alpha=\left(c,\frac{W_A-uc}{v}\right)\in\Newt(A),
 \qquad
 \beta=\left(c,\frac{W_B-uc}{v}\right)\in\Newt(B).
\]
Their difference is
$\delta=(0,(W_B-W_A)/v)\in\Q_{\geq0}^2$.  Every coefficient of
$A(\partial)^m(B^m)$ vanishes, including the coefficient on the ray
$m\delta$.  Proposition~\ref{prop:shifted} contradicts
$\alpha\in\Newt(A)\cap(\Newt(B)-\delta)$.
\end{proof}

Taking $W_A=W_B$ and $\delta=0$ in
Proposition~\ref{prop:shifted} gives the terminal version: equal-weight
pure-zero faces have disjoint Newton segments.  Because they lie on the
same line, one coordinate has a strict integral gap between all operator
exponents and all polynomial exponents.

\section{A common threshold is terminal}

The preceding coordinate gap survives bounded excursions away from the
two faces.

\begin{proposition}[Common-threshold cutoff]\label{prop:threshold}
Let $w=(u,v)$ be a positive unequal integral weight.  Suppose there is a
number $W$ such that
\begin{equation}\label{eq:threshold}
 w(\alpha)\geq W\quad(\alpha\in\Supp\lambda),
 \qquad
 w(\beta)\leq W\quad(\beta\in\Supp P).
\end{equation}
If $\Lambda^m(P^m)=0$ for every $m\geq1$, then
$\Lambda^m(QP^m)=0$ for every fixed $Q$ and all large $m$.
\end{proposition}

\begin{proof}
Let $A$ be the weight-$W$ face of $\lambda$ and $B$ the weight-$W$ face
of $P$.  The part of greatest possible output weight in the pure identity
is
\[
 A(\partial)^m(B^m)=0.
\]
By Proposition~\ref{prop:shifted} with $\delta=0$, the two Newton segments
are disjoint.  After possibly exchanging $x$ and $y$, there is an integer
$\epsilon>0$ such that
\begin{equation}\label{eq:coordinate-gap}
 \alpha_x\geq\beta_x+\epsilon
 \qquad
 (\alpha\in\Supp A,\ \beta\in\Supp B).
\end{equation}

Consider one monomial selection in $\Lambda^m(QP^m)$.  For an operator
selection $\alpha_i$ and a polynomial selection $\beta_j$, define their
nonnegative integral weight defects
\[
 a_i=w(\alpha_i)-W,
 \qquad b_j=W-w(\beta_j).
\]
If the selected monomial of $Q$ is $q$, a nonzero output monomial must have
nonnegative weight.  Therefore
\[
 \sum_i a_i+\sum_j b_j\leq w(q)\leq
 K:=\max_{q\in\Supp Q}w(q).
\]
At most $K$ operator selections and at most $K$ polynomial selections can
lie off the equality faces.  The full supports are finite, so their total
effect on the $x$-coordinate is bounded independently of $m$.  The remaining
face selections contribute, by \eqref{eq:coordinate-gap}, an $x$-derivative
excess at least $m\epsilon-C_Q$.  For large $m$ this exceeds the
$x$-degree supplied by $Q$.  The selected derivative therefore cannot act
nontrivially.  Every selection vanishes, proving the proposition.
\end{proof}

\section{The moving envelopes}

We now prove Theorem~\ref{thm:main}.  For $s\geq1$, give $x,X$ weight $s$
and $y,Y$ weight $1$.  Define
\begin{equation}\label{eq:envelopes}
 L(s)=\min_{\alpha\in\Supp\lambda}(s\alpha_x+\alpha_y),
 \qquad
 U(s)=\max_{\beta\in\Supp P}(s\beta_x+\beta_y),
 \qquad
 \Delta(s)=U(s)-L(s).
\end{equation}
These are finite, continuous, piecewise-linear functions.

By Lemma~\ref{lem:hall}, the least $x$-exponent on the lowest ordinary
operator face is $e$, while the greatest $x$-exponent on the top polynomial
face is some $t\leq e-1$.  For every sufficiently small rational
$\varepsilon>0$, the extremal faces at $s=1+\varepsilon$ are the
corresponding endpoints.  Since $d>r$,
\begin{equation}\label{eq:initial}
 \Delta(1+\varepsilon)>0,
 \qquad e>t.
\end{equation}

At a rational $s>1$, let $A_s$ be the complete $L(s)$-face of $\lambda$
and let $B_s$ be the complete $U(s)$-face of $P$.  The part of largest
$(s,1)$-weight in the pure identity is
\begin{equation}\label{eq:face-pure}
 A_s(\partial)^m(B_s^m)=0\qquad(m\geq1).
\end{equation}
Whenever $\Delta(s)>0$, Corollary~\ref{cor:horizontal} says that the two
$x$-intervals are disjoint.

Fix a sufficiently small rational $\varepsilon>0$ for which
\eqref{eq:initial} holds.  We prove that $\Delta$ reaches zero on the
connected interval beginning at $1+\varepsilon$.  Suppose, to the contrary,
that
\begin{equation}\label{eq:positive-forever}
 \Delta(s)>0\qquad(s\geq1+\varepsilon).
\end{equation}
As $s$ increases, the active slope of the lower envelope $L$ is
non-increasing, while the active slope of the upper envelope $U$ is
non-decreasing.  Equivalently, the active operator $x$-exponents move only
to the left and the active polynomial $x$-exponents move only to the right.
They begin with the complete operator interval strictly to the right of the
complete polynomial interval by \eqref{eq:initial}.

This order cannot reverse under \eqref{eq:positive-forever}.  A first
reversal could occur only at one of the finitely many rational breakpoints.
At that breakpoint, the complete exposed face of each envelope contains
its active exponents on both adjacent linear pieces.  The monotonicity just
noted then forces the two complete $x$-intervals to overlap.  Since the
breakpoint is rational and $\Delta>0$ there, this contradicts
Corollary~\ref{cor:horizontal}.  Consequently the operator interval remains
strictly to the right for every $s\geq1+\varepsilon$.

After the last breakpoint, $L$ is realized by an operator monomial with
globally least $X$-exponent $a_{\min}$, while $U$ is realized by a polynomial
monomial with globally greatest $x$-exponent $i_{\max}$.  The preserved
order gives $a_{\min}>i_{\max}$.  Hence on the final linear interval
\[
 \Delta(s)=C+s(i_{\max}-a_{\min})
\]
has strictly negative slope and must cross zero.  This contradicts
\eqref{eq:positive-forever}.

By continuity, there is therefore a first
$s_*>1+\varepsilon$ with $\Delta(s_*)=0$ on this initial positive
component.  All breakpoints and zeroes of these
integer-exponent envelopes are rational.  After multiplying the weight
$(s_*,1)$ by a common denominator, we obtain a positive unequal integral
weight.  With $W=L(s_*)=U(s_*)$, the definitions give
\[
 w(\alpha)\geq W\quad(\alpha\in\Supp\lambda),
 \qquad
 w(\beta)\leq W\quad(\beta\in\Supp P).
\]
Proposition~\ref{prop:threshold} proves eventual mixed vanishing.  Together
with Proposition~\ref{prop:cutoff} and the trivial constant-term cases,
this completes the proof of Theorem~\ref{thm:main}.  \qed

\section{From Long's Gaussian example to GVC}

Long's three-variable Gaussian counterexample is most transparent in
circular coordinates
\[
 Z=\frac{X+iY}{\sqrt2},\qquad
 W=\frac{X-iY}{\sqrt2},
\]
with $T$ an independent standard real Gaussian.  He takes
\begin{equation}\label{eq:long}
 P_{\mathrm L}=(1+Z)\left(W-\frac{2+Z}{2}T^2\right),
 \qquad Q_{\mathrm L}=Z,
\end{equation}
and proves the stronger coefficient identity
\begin{equation}\label{eq:long-coefficient}
 \mathbb E\bigl(A(Z)P_{\mathrm L}^m\bigr)
 =m![z^m]A(z)(1+z)^{m-1}.
\end{equation}
Thus the pure moment, obtained from $A=1$, is zero, while $A=z$ gives
$m!$~\cite{LongGMC}.

The feature we need is the adjacent-coefficient failure in
\eqref{eq:long-coefficient}: the pure functional asks one step beyond the
degree of $(1+z)^{m-1}$, while multiplication by $Z$ shifts the request
back to its leading coefficient.  A Gaussian expectation is not literally
an iterate of one fixed differential operator on powers of one homogeneous
polynomial, so an additional step is required for GVC.  We preserve the
same adjacent-coefficient mechanism on a quadric, increase the endpoint
contact from one to two, and thereby make the Laurent expression
polynomial and homogeneous in three variables.

\section{A homogeneous counterexample in three variables}

Work first over $\mathbb C$.  In variables $x,y,t$, put
\begin{equation}\label{eq:gvc3-data}
 \rho=t^2+xy,
 \qquad A=\rho+x^2,
 \qquad C=y\rho^2-2xt^2\rho-x^3t^2,
\end{equation}
and define
\begin{equation}\label{eq:gvc3-pair}
 P=AC^2,
 \qquad
 \Delta=4\partial_x\partial_y+\partial_t^2,
 \qquad
 \Lambda=\Delta^6,
 \qquad Q=x^2.
\end{equation}
The polynomial $P$ is homogeneous of degree twelve and $\Lambda$ is
homogeneous of order twelve.

\begin{theorem}[Three-variable counterexample]\label{thm:gvc3}
For every $m\geq1$,
\begin{equation}\label{eq:gvc3-pure-mixed}
 \Lambda^m(P^m)=0,
 \qquad
 \Lambda^m(x^2P^m)\ne0.
\end{equation}
More precisely,
\begin{equation}\label{eq:gvc3-exact}
 \Delta^{6m+1}(x^2P^m)
 =2^{8m+1}(6m+1)!(2m)!
   \frac{(12m+3)!!}{(4m+1)!!}.
\end{equation}
The same identities hold over every characteristic-zero field.
\end{theorem}

\subsection{The polynomial lift}

The reason for the definition of $C$ is the identity
\begin{equation}\label{eq:cusp-identity}
 xC=\rho^3-t^2(\rho+x^2)^2.
\end{equation}
Indeed, $xy=\rho-t^2$, so
\[
 xC=(\rho-t^2)\rho^2-2x^2t^2\rho-x^4t^2
    =\rho^3-t^2(\rho+x^2)^2.
\]
On the affine quadric $\rho=1$, formula \eqref{eq:cusp-identity} gives
\begin{equation}\label{eq:quadric-laurent}
 P=x^{-2}(1+x^2)
       \left(1-t^2(1+x^2)^2\right)^2.
\end{equation}
This Laurent expression is used only for coefficient extraction;
\eqref{eq:gvc3-data}--\eqref{eq:gvc3-pair} show that $P$ itself is a
polynomial.  The square is the extra endpoint contact which removes the
pole after homogenization.

\subsection{The two neighboring spherical coefficients}

Write $x=X+iY$, $y=X-iY$, and keep $t=T$.  Then
\[
 \rho=X^2+Y^2+T^2,
 \qquad
 \Delta=\partial_X^2+\partial_Y^2+\partial_T^2.
\]
On the unit sphere, parameterize
\[
 x=\sqrt{1-t^2}\,e^{i\theta},
 \qquad y=\sqrt{1-t^2}\,e^{-i\theta},
\]
with normalized measure $(dt/2)(d\theta/2\pi)$.  Phase integration in
\eqref{eq:quadric-laurent} is constant-term extraction.  Setting $u=x^2$
gives
\begin{equation}\label{eq:sphere-coefficients}
 \int_{S^2}P^m\,d\sigma=[u^m]K_m(u),
 \qquad
 \int_{S^2}x^2P^m\,d\sigma=[u^{m-1}]K_m(u),
\end{equation}
where
\begin{equation}\label{eq:km}
 K_m(u)=(1+u)^m\int_0^1
       \left(1-v^2(1+u)^2\right)^{2m}\,dv.
\end{equation}

Put $B=1+u$ and
\[
 J_m(B)=\int_0^B(1-w^2)^{2m}\,dw.
\]
Changing variables in \eqref{eq:km} gives
\[
 K_m(u)=B^{m-1}J_m(B).
\]
Because $J_m'(B)=(1-B^2)^{2m}$ has a zero of order $2m$ at $B=1$,
\[
 J_m(1+u)=J_m(1)+O(u^{2m+1}).
\]
Consequently, through degree $m$,
\[
 K_m(u)=J_m(1)(1+u)^{m-1}+O(u^{m+1}).
\]
The coefficient of $u^m$ is zero and the coefficient of $u^{m-1}$ is
\begin{equation}\label{eq:cm}
 C_m=J_m(1)=\int_0^1(1-w^2)^{2m}\,dw
 =\frac{2^{2m}(2m)!}{(4m+1)!!}\ne0.
\end{equation}
This is exactly the pure-versus-adjacent coefficient split seen in Long's
identity \eqref{eq:long-coefficient}.

\subsection{From sphere averages to derivatives}

For every homogeneous polynomial $F$ of degree $2k$ in three variables,
orthogonal invariance gives the apolar Reynolds identity
\begin{equation}\label{eq:reynolds}
 \Delta^kF
 =2^kk!(2k+1)!!\int_{S^2}F\,d\sigma.
\end{equation}
Both sides are invariant linear functionals on homogeneous degree-$2k$
polynomials; evaluation on $\rho^k$ fixes the displayed constant, so
\eqref{eq:reynolds} is algebraic despite the convenient spherical notation.

Since $P^m$ has degree $12m$, equations
\eqref{eq:sphere-coefficients}, \eqref{eq:cm}, and
\eqref{eq:reynolds} give
\[
 \Delta^{6m}(P^m)=0.
\]
Since $x^2P^m$ has degree $12m+2$, they also give
\begin{align*}
 \Delta^{6m+1}(x^2P^m)
 &=2^{6m+1}(6m+1)!(12m+3)!!\,C_m\\
 &=2^{8m+1}(6m+1)!(2m)!
   \frac{(12m+3)!!}{(4m+1)!!},
\end{align*}
which is nonzero.  If $\Delta^{6m}(x^2P^m)$ were zero, applying one more
$\Delta$ would give zero, so the mixed output in
\eqref{eq:gvc3-pure-mixed} is nonzero.  All displayed differential
identities have rational coefficients.  The result therefore descends from
$\mathbb C$ to every characteristic-zero field, proving
Theorem~\ref{thm:gvc3}.

\section{The full cusp-profile family}

The construction is determined by the coefficient calculation.
Long's example suggested looking for a phase variable $u$ for which the
pure moment asks for the first coefficient beyond
$(1+u)^{m-1}$, while multiplication by $u$ moves the request back to the
last nonzero coefficient.  On the quadric $\rho=1$, the choice
\[
 B=1+x^2,
 \qquad E=1-t^2B^2
\]
has the required endpoint flatness at $B=1$.  The desired phase winding is
$x^{-2}$.  A single $E$ has only one hidden factor of $x$, so it does not
remove that pole.  The identity
\[
 E=xC
\]
on $\rho=1$ shows that $E^2x^{-2}=C^2$ is the first polynomial choice.
Multiplying by $A=\rho+x^2$ supplies the binomial $B$.  This forces
\[
 P=AC^2
\]
and explains both the square and degree twelve: they are not guesses from
an expanded 23-term polynomial, but the smallest contact-two homogeneous
lift in this one-profile construction.

The same design has independent winding, profile, and radial parameters.
Let $r\geq1$, let $e,h\geq0$, and choose
\[
 S(z)=\sum_{j=0}^e s_jz^j\ne0.
\]
Define its homogeneous lift and the resulting polynomial by
\begin{equation}\label{eq:profile-family}
 \begin{aligned}
 S^{\mathrm{hom}}
   &=\sum_{j=0}^e
      s_j(t^2A^2)^j\rho^{3(e-j)},\\
 P_{r,S,h}
   &=\rho^hA^rC^{2r}S^{\mathrm{hom}},\\
 N&=6r+3e+h,
 \qquad \Lambda_{r,S,h}=\Delta^N.
 \end{aligned}
\end{equation}
Then $P_{r,S,h}$ is homogeneous of degree $2N$.  Put
\begin{equation}\label{eq:profile-moment}
 c_m(S)=\int_0^1
   (1-v^2)^{2rm}S(v^2)^m\,dv.
\end{equation}
Algebraically, if the integrand is $\sum_j a_jv^{2j}$, this notation means
$\sum_j a_j/(2j+1)$; it therefore makes sense over every
characteristic-zero field.

\begin{theorem}[Cusp-profile failure family]\label{thm:family}
If $c_m(S)\ne0$ for every $m\geq1$, then
\begin{equation}\label{eq:family-pure}
 \Delta^{Nm}(P_{r,S,h}^m)=0
 \qquad(m\geq1),
\end{equation}
while, for $1\leq\ell\leq rm$,
\begin{equation}\label{eq:family-ladder}
 \begin{aligned}
 &\Delta^{Nm+\ell}
   \left(x^{2\ell}P_{r,S,h}^m\right)\\
 &\quad={}
 2^{Nm+\ell}(Nm+\ell)!(2Nm+2\ell+1)!!
 \binom{rm-1}{\ell-1}c_m(S)\ne0.
 \end{aligned}
\end{equation}
In particular, $Q=x^2$ disproves GVC for every member of the family at
every positive power.
\end{theorem}

\begin{proof}
On $\rho=1$, put $u=x^2$ and $B=1+u$.  Formula
\eqref{eq:profile-family} becomes
\[
 \left.P_{r,S,h}\right|_{\rho=1}
 =u^{-r}B^r(1-t^2B^2)^{2r}S(t^2B^2).
\]
After the $m$-th power, phase extraction asks for the coefficients
$u^{rm}$ and $u^{rm-\ell}$ of
\[
 H_m(B)=B^{rm-1}J_m(B),
 \qquad
 J_m'(B)=\left((1-B^2)^{2r}S(B^2)\right)^m.
\]
The derivative has a zero of order at least $2rm$ at $B=1$, so
\[
 J_m(1+u)=c_m(S)+O(u^{2rm+1}).
\]
Through degree $rm$, the relevant polynomial is therefore
$c_m(S)(1+u)^{rm-1}$.  Its $u^{rm}$ coefficient is zero and its
$u^{rm-\ell}$ coefficient is
$\binom{rm-1}{\ell-1}c_m(S)$.  Applying the Reynolds identity
\eqref{eq:reynolds} proves \eqref{eq:family-pure} and
\eqref{eq:family-ladder}.
\end{proof}

The condition in Theorem~\ref{thm:family} is easy to satisfy.  For example,
$S(z)>0$ on $[0,1]$ gives $c_m(S)>0$ for every $m$, so each fixed
$(r,e,h)$ contains an open real coefficient family rather than one isolated
point.  Three useful specializations are:
\begin{itemize}
 \item $r=1$, $S=1$, $h=0$: the minimal witness $AC^2$;
 \item for every $s\geq2$, take $S(z)=(1-z)^{s-2}$ and $h=0$:
 \[
   P_s=Ax^{s-2}C^s,
   \qquad \Lambda_s=\Delta^{3s};
 \]
 \item for every $k\geq6$, take $r=1$, $S=1$, and $h=k-6$:
 \begin{equation}\label{eq:radial-family}
   P_k=\rho^{k-6}AC^2,
   \qquad \Lambda_k=\Delta^k.
 \end{equation}
\end{itemize}
Thus every power $\Delta^k$ with $k\geq6$ has a homogeneous
three-variable GVC counterexample, and every such example carries the full
positive-phase multiplier ladder in \eqref{eq:family-ladder}.

\section{The exact dimensional boundary}

Let $\operatorname{GVC}(n)$ denote \eqref{eq:gvc} for all
constant-coefficient operators and all polynomials in $n$ variables, and
let $\operatorname{HGVC}(n)$ denote the same assertion restricted to
homogeneous operators.

\begin{lemma}[Unused-variable padding]\label{lem:padding}
Let $1\leq r\leq n$.  Every counterexample to $\operatorname{GVC}(r)$
over a characteristic-zero field gives a counterexample to
$\operatorname{GVC}(n)$ over the same field.  The same assertion holds for
homogeneous GVC.
\end{lemma}

\begin{proof}
Write the counterexample in variables $x_1,\ldots,x_r$ as
$\Lambda=\lambda(\partial_{x_1},\ldots,\partial_{x_r})$, $P$, and $Q$.
View $\lambda$, $P$, and $Q$ as polynomials in
$x_1,\ldots,x_n$ that are independent of the added variables
$x_{r+1},\ldots,x_n$, and let $\widetilde\Lambda$ be the resulting
$n$-variable operator.  If
\[
 \iota:k[x_1,\ldots,x_r]\longrightarrow k[x_1,\ldots,x_n]
\]
is the natural inclusion, then coefficientwise differentiation gives, for
every $m\geq1$,
\[
 \widetilde\Lambda^m\bigl(\iota(P)^m\bigr)
 =\iota\bigl(\Lambda^m(P^m)\bigr),
 \qquad
 \widetilde\Lambda^m\bigl(\iota(Q)\iota(P)^m\bigr)
 =\iota\bigl(\Lambda^m(QP^m)\bigr).
\]
Thus all pure identities are preserved, as is every nonzero mixed output.
The inclusion also preserves total degrees, so a homogeneous operator and
homogeneous polynomial remain homogeneous after padding.
\end{proof}

\begin{theorem}[Dimension classification]\label{thm:classification}
For every $n\geq1$,
\begin{equation}\label{eq:classification}
 \operatorname{GVC}(n)\text{ holds}
 \quad\Longleftrightarrow\quad
 n\leq2.
\end{equation}
The same equivalence holds for $\operatorname{HGVC}(n)$.
\end{theorem}

\begin{proof}
Theorem~\ref{thm:main} proves the unrestricted assertion for two variables,
and the one-variable case follows by restriction or by the split-symbol
theorem.  For the homogeneous assertion in at most two variables, every
homogeneous symbol splits into linear forms after scalar extension, which
is the case proved by de Bondt~\cite{deBondt}.

Theorem~\ref{thm:gvc3} gives a homogeneous failure in three variables.
Lemma~\ref{lem:padding} transports this same counterexample to every
$n\geq3$, while preserving homogeneity.  Equivalently,
$\neg\operatorname{GVC}(3)$ implies
$\neg\operatorname{GVC}(n)$ and
$\neg\operatorname{HGVC}(n)$ for every $n\geq3$.
\end{proof}

The classification is a closure of GVC by dimension, not a solution of the
Jacobian Conjecture.  Nor does Theorem~\ref{thm:gvc3} settle Zhao's
ordinary-Laplacian Vanishing Conjecture: its operator is the sixth power
$\Delta^6$, not $\Delta$ itself.

\section{Why the two-dimensional proof stops}

The positive and negative parts of the paper meet at a structural fault
line.  In two variables, every homogeneous symbol splits after scalar
extension, Hall deficiency localizes at one repeated direction, and an
unequal-weight Newton face is a segment.  Those facts force the moving
operator and polynomial intervals to meet at a terminal threshold.

The three-variable symbol
\[
 (4XY+T^2)^6
\]
contains a rank-three quadratic factor and does not split into linear forms.
On the sphere, its invariant functional retains the radial variable while
phase integration extracts one coefficient.  The construction arranges
for the pure coefficient to lie just beyond the degree of
$(1+u)^{m-1}$, but multiplication by $x^2$ shifts the extraction to the
nonzero neighboring coefficient.  There is no single horizontal interval
gap for the binary envelope argument to preserve.  This is the precise
mechanism of failure, rather than merely a dimension count.

\section{A Ferrers support consequence}

Suppose a wall begins with a horizontal gap $g>0$.
Index tied operator corrections by their leftward moves $i$ and polynomial
corrections by their rightward moves $j$.  The two face intervals overlap
exactly when
\[
 i+j\geq g.
\]
Horizontal separation forbids every such pair.  Set-theoretically, the
allowed supports are therefore the independent sets of the Ferrers graph
with edges $i+j\geq g$.  The corresponding squarefree ideal is
\[
 (a_i:i\geq g)+(b_j:j\geq g)
 +(a_i b_j:i,j<g,\ i+j\geq g).
\]
Thus the squarefree support obstruction has coordinate components and
staircase products as a direct consequence of horizontal separation.

The proof is specifically two-dimensional.  Hall deficiency localizes at
one repeated binary direction, and an unequal-weight face is a segment
whose projection is an interval.  In higher dimension, a deficient set can
contain a more complicated arrangement of directions, and two disjoint
faces can move around one another without passing through one common
one-dimensional projection.  The binary separation theorem therefore does
not extend formally to the rank-three symbol in
Theorem~\ref{thm:gvc3}.

\paragraph{Machine-checked audit.}
The accompanying Lean development in \texttt{formal/gvc} verifies the
coefficientwise differential semantics; the literal degree-twelve ternary
data and cusp identity; the algebraic beta evaluation and complete
cusp-profile binomial ladder; the literal winding--profile--radial family,
its degree formula, and its minimal and radial specializations;
equal-degree coefficientwise apolar
contraction and composition of differential symbols; the full concrete
Reynolds/Laurent phase extraction, algebraic height change of variables, and
endpoint-kernel identification.  It also proves the arbitrary-profile
quadric restriction under the displayed degree bound on $S$, arbitrary
even-phase coefficient extraction, the full height substitution and
endpoint factorization, and the shifted primitive identity.  These results
construct both phase bridges in Lean and prove
Theorems~\ref{thm:gvc3} and \ref{thm:family}, including their exact mixed
scalars and coefficient base change to every characteristic-zero field.
Unused-variable padding verifies the negative dimension-$n\ge3$ half of
Theorem~\ref{thm:classification}.  The development also checks the rational
$p$-adic factorial bounds in
Lemma~\ref{lem:factorial-valuations} and the negative-final-slope plus
intermediate-value step for the moving envelopes.  For
Lemma~\ref{lem:hall}, it proves the finite Hall-deficiency core: Mathlib's
Hall theorem produces a deficient set, binary linear algebra puts that set
in one direction class, and exact counting gives at least $d-e+1$
annihilated polynomial factors.  The split-factor interface records the
literal symbol and polynomial products, and Lean proves the coordinate-free
normal forms: the full direction power divides the symbol and the matching
perpendicular power divides the polynomial.  The translated
Duistermaat--van der Kallen/polarization argument that rules out a matching
is not yet formalized.  It further proves the
complete finite-support deduction in Proposition~\ref{prop:threshold}:
finite maximization turns an ordering of the equal faces into an integral
coordinate cut, which refines the common threshold to a strict positive
weight; the growing power gap then dominates every fixed multiplier.  Empty
equality faces are handled separately.
The present audit does not yet formalize construction of the split-factor
system after field extension or the literal coordinate-change/exact-quotient
packaging of Lemma~\ref{lem:hall}, the number-field
good-prime transfer and shifted-ray theorem, including its equal-face
ordering, or the binary no-reversal/global-envelope bridge.  The spherical realization of the
Reynolds functional is likewise not checked, but Lean proves the algebraic
coefficientwise contraction and Laurent realization used in its place.
Thus Theorem~\ref{thm:main} and the complete biconditional in
Theorem~\ref{thm:classification} are not yet fully formally verified.

\bibliographystyle{plain}
\bibliography{references}

\end{document}
